\section{结论}

本文围绕霍尔木兹海峡封锁对国际原油价格的影响，建立传统供需基准模型、短期冲击模型和中长期油价调节模型，得到以下结论。

第一，在赛题低短期弹性设定下，静态反事实基准给出 278--337 美元/桶的压力上界。该数值刻画低弹性框架下的边际压力，不代表现实价格必然到达。其作用是说明：若忽略 SPR、商业库存、绕道运输、需求收缩和市场预期修复，传统供需框架会系统性高估现实价格。

第二，短期冲击模型能够解释冲突窗口内约 110--120 美元/桶的价格平台。加入 SPR、商业库存、绕道运输、需求收缩、恐慌衰减、地缘风险溢价和预期修复后，模型在冲突窗口内取得 RMSE 3.47 美元/桶、MAE 2.87 美元/桶，并优于朴素上一日基准和滚动 ARIMA 基准。短期平台并非简单复制昨日价格，而是物理缺口、政策缓冲和市场预期共同再定价的结果。

第三，中长期油价调节模型显示，中性情景第 180 天价格约为 97.82 美元/桶；悲观情景在供应中断偏大、SPR 等效预算耗尽、绕道恢复缓慢和制度风险溢价持续时，第 180 天价格约为 121.64 美元/桶，并保留二次跳涨风险。敏感性分析表明，封锁风险衰减速度、SPR 释放上限、供应中断量和绕道运输能力是长期结果的主要驱动因素。

第四，外部风险变量用于约束风险环境而非替代附件价格。GPR 显示 2026 年 3--4 月地缘风险处于历史高分位，OVX 显示冲突窗口隐含波动率均值处于同长度历史窗口约 97.3\% 分位；2017--2026 年同长度窗口也表明，本题冲突窗口在累计收益率、波动率和朴素基准误差上均具有显著特殊性。模型适用于霍尔木兹封锁条件路径分析，不宜解释为跨年份通用短线交易器。

后续若能取得 IEA MODS、油轮保险费、BDTI/TD3C 和期货期限结构等授权数据，可进一步细化库存消耗、运输风险和时变状态转移概率。
